\uuid{3M43}
\exo7id{4442}
\titre{exo7 4442}
\auteur{quercia}
\organisation{exo7}
\datecreate{2010-03-14}
\isIndication{false}
\isCorrection{true}
\chapitre{Série numérique}
\sousChapitre{Autre}
\cpge{mpsi-mp2i.procedes-sommatoires-discrets.a, pcsi.series-numeriques.a, ptsi.series-numeriques.a, mpsi-mp2i.analyse-asymptotique.c, pcsi.analyse-asymptotique.c, ptsi.analyse-asymptotique.c, mpsi-mp2i.procedes-sommatoires-discrets.d, bcpst1.suites-reelles, mp-mpi.series-numeriques-vectorielles.b, pc.series-numeriques, psi.series-numeriques.a, pt.series-numeriques.a, bcpst2.series-reelles}
\module{Analyse}
\niveau{L1}
\difficulte{}

\contenu{
\question{
On pose $u_n = \sum_{k=1}^{2n} (-1)^k\sqrt k$. Donner un équivalent de $u_n$
quand $n\to\infty$.
(Regrouper les termes deux par deux puis comparer à une intégrale)
}
\reponse{
$\sqrt{\frac n2}$.
}
}
